\section{The first inverse-index correction}\label{sec:correction}

The leading amplitude only sees one large forward index and the linear
indices accompanying it. At the next order, a single quadratic forward
index also contributes. We retain that contribution uniformly in the number
of factors, then compute the corrections from the gamma ratios and from
$T_{n,k}$. This proves the full statement of Theorem~\ref{thm:main}.

\subsection{Two endpoint terms with a uniform remainder}

Recall the scaled coefficients and convolution coefficients from
Section~\ref{sec:leading}:
\[
 a_m=A^m(-1)^m r_m,\qquad
 S_{n,k}=[x^n]\left(\sum_{m\ge1}a_mx^m\right)^k.
\]
In particular $a_1=c$, $a_2=d$, and every $a_m$ is nonnegative. It is
convenient to shift the index once:
\begin{equation}\label{eq:alpha-shift}
 \alpha_j=a_{j+1},\qquad
 \alpha(x)=\sum_{j\ge0}\alpha_jx^j,\qquad
 S_{n,k}=[x^{n-k}]\alpha(x)^k.
\end{equation}
Thus $\alpha_0=c$, $\alpha_1=d$, and there is a constant $B$ such that
\begin{equation}\label{eq:alpha-bound}
 0\le\alpha_j\le B\Gamma(j)\qquad(j\ge1).
\end{equation}
Indeed $a_{j+1}=4\Gamma(j+1)\zeta(j+1)^2/(j+1)$ when $j+1$ is even,
and it is zero otherwise. The ratio $j/(j+1)$ is at most one, and
$\zeta(j+1)\le\zeta(2)$.

\begin{lemma}[The second endpoint contribution]\label{lem:second-endpoint}
There is a constant $B_*>0$ such that, whenever $n\ge12$ and
$1\le k\le n/2$,
\begin{equation}\label{eq:two-endpoints}
 \begin{split}
 S_{n,k}={}&kc^{k-1}a_{n-k+1}
       +k(k-1)c^{k-2}d\,a_{n-k}\\
       &\quad+O\bigl(B_*^k\Gamma(n-k-2)\bigr).
 \end{split}
\end{equation}
The second term is understood to be zero for $k=1$, and the implied
constant is independent of both $n$ and $k$.
\end{lemma}
\begin{proof}
We first record two elementary estimates. The sequence $F_j=j!$ satisfies
$F*F\le3F$ coefficientwise. For $M\ge1$, division of its $M$th convolution
coefficient by $M!$ gives
\[
 \sum_{j=0}^{M}\binom{M}{j}^{-1}\le3:
\]
the endpoints contribute two, and each of the $M-1$ interior terms is at
most $1/M$. The case $M=0$ is immediate. Iteration, followed by the change
$j_i\mapsto j_i-1$, gives
\begin{equation}\label{eq:positive-gamma-convolution}
 \sum_{\substack{j_1+\cdots+j_h=N\\j_i\ge1}}
       \prod_{i=1}^{h}\Gamma(j_i)
 \le3^{h-1}\Gamma(N-h+1),\qquad N\ge h\ge1.
\end{equation}

Removing the endpoints from the two-factor convolution gives a sharper
estimate:
\begin{equation}\label{eq:interior-gamma-convolution}
 \sum_{j=2}^{N-2}\Gamma(j)\Gamma(N-j)
       \le4\Gamma(N-2),\qquad N\ge4.
\end{equation}
For completeness, put $M=N-2$ and divide the left side by $M!$. The
result is $\sum_{i=1}^{M-1}\binom{M}{i}^{-1}$. For $M\ge3$, its two extreme
terms contribute $2/M$, and each of the remaining $M-3$ terms is at most
$2/[M(M-1)]$. The sum is at most $4/M$. The case $M=2$ also satisfies this
bound, proving \eqref{eq:interior-gamma-convolution}.

Now put $N=n-k$, so $N\ge k$ and $N\ge6$, and expand $\alpha(x)^k$
according to the number $h$ of positive indices. The case $h=1$ contributes
$kc^{k-1}\alpha_N$. For $h=2$, the indices $1,N-1$ and $N-1,1$ contribute
\[
 2\binom{k}{2}c^{k-2}\alpha_1\alpha_{N-1}
       =k(k-1)c^{k-2}d\alpha_{N-1}.
\]
By \eqref{eq:alpha-bound} and \eqref{eq:interior-gamma-convolution}, the
remaining $h=2$ terms are bounded by
$4\binom{k}{2}c^{k-2}B^2\Gamma(N-2)$.
For $h\ge3$, \eqref{eq:positive-gamma-convolution} bounds the contribution
by
\[
 \binom{k}{h}c^{k-h}B^h3^{h-1}\Gamma(N-h+1).
\]
Since $N\ge k\ge h$, one has
$\Gamma(N-h+1)\le\Gamma(N-2)$. The binomial theorem therefore bounds
the sum of all these remainders by $B_*^k\Gamma(N-2)$ after a fixed
increase of $B_*$. Substituting $\alpha_N=a_{n-k+1}$ and
$\alpha_{N-1}=a_{n-k}$ proves \eqref{eq:two-endpoints}.
\end{proof}

\subsection{The two elementary expansions}

The next estimates are uniform over the full range $k\le n/2$; their
constants may grow exponentially in $k$, because the Lagrange factor
$1/k!$ makes that growth summable. Define
\[
 Q_{n,k}(u)=n^{1-k}T_{n,k}(u),\qquad
 \beta_k=\frac{(k-1)(k+2)}2.
\]
For $|u|\le U$, the finite sum \eqref{eq:T} gives
\begin{equation}\label{eq:T-first-correction}
 Q_{n,k}(u)=1+\frac{k(k-1)u}{2n}
       +O_U\left(\frac{(1+k)^4e^{B_Uk}}{n^2}\right).
\end{equation}
To justify the uniform remainder, use
$\binom{k-1}{j}\le k^j/j!$ and $(k)_j\le(2k)^j$. The sum of the absolute
values of the terms with $j\ge2$ is at most
\[
 \sum_{j\ge2}\frac{(Uk^2/n)^j}{j!}
 \le\frac{U^2k^4}{2n^2}e^{Uk^2/n}.
\]
The inequality $k^2/n\le k/2$ proves
\eqref{eq:T-first-correction}. It also gives a bound
$|Q_{n,k}(u)|\le e^{Uk/2}$ that will be used for the endpoint remainder.

For the first endpoint, put $p=n-k+1$. When $p$ is even,
\begin{equation}\label{eq:gamma-first-correction}
 \frac{n^{k-1}a_p}{b_n}
 =4\left\{1+\frac{\beta_k}{n}
       +O\left(\frac{(1+k)^4e^{B_1k}}{n^2}\right)\right\};
\end{equation}
when $p$ is odd it is zero. The exact gamma ratio before multiplication
by $4\zeta(p)^2$ is
\begin{equation}\label{eq:gamma-ratio-exact}
 \frac{n^{k-1}b_{n-k+1}}{b_n}
 =\frac{n^k}{(n-k+1)\prod_{j=1}^{k-1}(n-j)}.
\end{equation}
Its logarithm equals $\beta_k/n+O(k^3/n^2)$: the offsets in the
denominator are $k-1,1,\ldots,k-1$, whose sum is $\beta_k$, and every
offset divided by $n$ is at most $1/2$. More explicitly,
$-\log(1-v)=v+O(v^2)$ uniformly for $0\le v\le1/2$. The logarithm
itself is $O(k^2/n)=O(k)$, so exponentiating introduces a remainder
$O((1+k)^4e^{B_1k}/n^2)$. Finally
$\zeta(p)^2=1+O(2^{-p})$, with $p\ge n/2$, and this exponentially small
error is absorbed by the displayed bound. This proves
\eqref{eq:gamma-first-correction}.

The second endpoint has the corresponding estimate
\begin{equation}\label{eq:gamma-second-endpoint}
 \frac{n^{k-1}a_{n-k}}{b_n}
 =\frac4n+O\left(\frac{(1+k)^2e^{B_2k}}{n^2}\right)
       \quad\text{if }n-k\text{ is even},
\end{equation}
and it is zero otherwise. Here the exact gamma ratio is
\[
 \frac{n^{k-1}b_{n-k}}{b_n}
 =\frac{n^k}{(n-k)\prod_{j=1}^{k}(n-j)}.
\]
Factor out $1/n$ and apply the same logarithmic estimate to the offsets
$k,1,\ldots,k$. This proves \eqref{eq:gamma-second-endpoint}.

The remainder in Lemma~\ref{lem:second-endpoint} is also uniformly of
second order after normalization. For $n\ge12$ and $k\le n/2$,
\begin{equation}\label{eq:endpoint-normalized-error}
 \frac{n^{k-1}\Gamma(n-k-2)}{b_n}
 =\frac{n^k\Gamma(n-k-2)}{\Gamma(n)}
 \le\frac{3^{k+2}}{n^2}.
\end{equation}
Indeed, the denominator ratio contains $k+2$ factors $n-j$, each at least
$n-k-2\ge n/3$. Combining
\eqref{eq:two-endpoints}--\eqref{eq:endpoint-normalized-error}, all errors
in the $k$th summand of \eqref{eq:C-finite}, after division by $b_n$ and
scaling by $A^n$, are bounded by
\begin{equation}\label{eq:summable-second-error}
 \frac{|u|}{n^2}\frac{M_U^k(1+k)^M}{k!},
\end{equation}
for constants $M_U$ and a fixed integer $M$. One obtains this bound by
leaving one factor $|u|$ outside and replacing the remaining
$|u|^{k-1}$ by $U^{k-1}$. Thus the estimates remain uniform after division
by $u$, with the removable value at zero.

\subsection{Summing the corrections}

We now complete the proof of Theorem~\ref{thm:main}. Recall that
$D_n(u)=(-1)^{n+1}A^nC_n(u)$. The range $k>n/2$ contributes less than
any prescribed inverse power of $n$ after normalization. To make the
uniformity at $u=0$ explicit, the estimate
\eqref{eq:master-majorant}, with one factor $u$ left outside, gives
\begin{equation}\label{eq:correction-tail-majorant}
 \left|\frac{u^kS_{n,k}T_{n,k}(u)}{k!b_n}\right|
 \le |u|L_U\frac{kL_U^{k-1}}{k!},\qquad 1\le k\le n.
\end{equation}
It follows equally by inserting the three bounds in
Lemma~\ref{lem:factorial} into the left side. For every fixed integer
$J\ge1$, use $1\le(2k/n)^J$ when $k>n/2$ to obtain
\[
 \sum_{k>n/2}
 \left|\frac{u^kS_{n,k}T_{n,k}(u)}{k!b_n}\right|
 \le\frac{|u|\,2^JL_U}{n^J}
       \sum_{k\ge1}\frac{k^{J+1}L_U^{k-1}}{k!}
 =O_{U,J}(|u|n^{-J}).
\]
The tails of the exponential series below, including their polynomial
weights in $k$, satisfy the same estimate. This permits replacement of
the remaining finite sums by infinite sums at the required error order.

For $k\le n/2$, the first endpoint in \eqref{eq:two-endpoints} contributes
only when $n-k+1$ is even. Put $\ell=k-1$; its parity condition becomes
$\ell\equiv n\pmod2$. After division by $4ub_n$, its contribution is
\[
 \frac{z^\ell}{\ell!}
 \left[1+\frac{\ell(\ell+3)+u\ell(\ell+1)}{2n}\right],
 \qquad z=cu,
\]
up to the summable error \eqref{eq:summable-second-error} divided by
$4|u|$. For the second endpoint put $\ell=k-2$. Its parity condition is
again $\ell\equiv n\pmod2$, and its first-order contribution is
$du\,z^\ell/(n\ell!)$. Hence
\begin{equation}\label{eq:correction-parity-sums}
 \begin{split}
 \frac{D_n(u)}{4ub_n}
 ={}&\sum_{\substack{\ell\ge0\\\ell\equiv n\ (2)}}\frac{z^\ell}{\ell!}\\
 &+\frac1n\sum_{\substack{\ell\ge0\\\ell\equiv n\ (2)}}
 \left\{du+\frac{\ell(\ell+3)+u\ell(\ell+1)}2\right\}
 \frac{z^\ell}{\ell!}+O_U(n^{-2}).
 \end{split}
\end{equation}
All sums and remainders are locally uniform in complex $u$.

If $H=H_{n\bmod2}$, the elementary identities
\[
 \sum_{\ell\equiv n\ (2)}\ell(\ell+3)\frac{z^\ell}{\ell!}
       =z^2H(z)+4zH'(z),
 \qquad
 \sum_{\ell\equiv n\ (2)}\ell(\ell+1)\frac{z^\ell}{\ell!}
       =z^2H(z)+2zH'(z)
\]
turn \eqref{eq:correction-parity-sums} into
\eqref{eq:main-expansion}--\eqref{eq:K}.

It remains to prove the extra factor $u$ in the odd remainder. For odd
$n\ge3$, let
\[
 E_n(u)=\frac{D_n(u)}{4ub_n}-H_1(cu)-\frac{K_1(u)}n.
\]
This is entire in $u$. By the exact valuation in \eqref{eq:valuation},
and by the displayed formulas for $H_1$ and $K_1$, it vanishes at zero.
The already proved estimate on the larger disk $|u|\le2U$ bounds
$E_n$ there by $O_{2U}(n^{-2})$. Apply the maximum-modulus principle to
the entire function $E_n(u)/u$ on that disk: its boundary values are
$O_{2U}(n^{-2}/U)$. Consequently
\begin{equation}\label{eq:odd-extra-u}
 E_n(u)=O_U(|u|n^{-2}),\qquad |u|\le U.
\end{equation}
This completes the complex-uniform part of Theorem~\ref{thm:main}.

\begin{corollary}[Relative uniformity down to the resonant endpoint]
\label{cor:relative-two-term}
For every fixed $U>0$, uniformly for real $0<u\le U$,
\begin{equation}\label{eq:relative-two-term}
 C_n(u)=(-1)^{n+1}\frac{4ub_n}{A^n}H_{n\bmod2}(cu)
 \left\{1+\frac{K_{n\bmod2}(u)}{nH_{n\bmod2}(cu)}
                +O_U(n^{-2})\right\}.
\end{equation}
The ratios $K_0(u)/H_0(cu)$ and $K_1(u)/H_1(cu)$ extend continuously to
$u=0$, with respective limits $0$ and $2$.
\end{corollary}
\begin{proof}
For $u\ge0$, $\cosh(cu)\ge1$ and $\sinh(cu)\ge cu$. Divide the even
remainder by the former bound and the sharpened odd remainder
\eqref{eq:odd-extra-u} by the latter. Formula~\eqref{eq:K} gives the
stated limits. In particular the ratios are bounded on $[0,U]$, and
\eqref{eq:even-main}--\eqref{eq:odd-main} follow as well.
\end{proof}

\subsection{An exact check at small slow parameter}

The quadratic coefficient in Proposition~\ref{prop:structure} provides an
all-index consistency check. For odd $n\ge3$,
\begin{equation}\label{eq:small-u-exact}
 \begin{split}
 \lim_{u\to0}\frac{A^n C_n(u)}{4u^2b_n}
 &=\frac{A^n}{4b_n}\frac n4r_{n-1}\\
 &=c\frac{n^2}{(n-1)^2}\zeta(n-1)^2
 =c\left\{1+\frac2n+O(n^{-2})\right\}.
 \end{split}
\end{equation}
Our expansion gives the same result: $H_1(cu)/u\to c$ and
$K_1(u)/u\to2c$, with a remainder that stays uniform after this division.
For even $n$, the corresponding linear-coefficient identity is
\[
 \lim_{u\to0}\frac{-A^nC_n(u)}{4ub_n}=\zeta(n)^2
       =1+O(2^{-n}).
\]
These checks also explain the different powers of $u$ in the two parity
sectors. If $u=u_n>0$ tends to zero at any rate while $n$ tends to infinity,
\eqref{eq:relative-two-term} gives
\[
 C_n(u_n)\sim
 \begin{cases}
 -4u_n b_nA^{-n},&n\text{ even},\\
 A u_n^2b_nA^{-n},&n\text{ odd}.
 \end{cases}
\]
Thus the late-coefficient theorem remains informative when the slow
parameter follows the physical specialization $u=X^{-1}$.
